\documentclass[9pt,aspectratio=169]{beamer}
\usetheme{StAndrews}

\coursename{FI5699 Dissertation Module}
\graphicspath{ {./logos/} {./outputs/} {./pix/week2/} }
\title{\LARGE \textnormal{Week 2: \\
Introduction to Python}}
\hypersetup{pdftitle={FI5699 Week 2 -- Introduction to Python}}

% ============================================================================
% CODE LISTINGS SETUP
% ============================================================================
\usepackage{listings}
\lstset{
  language=Python,
  basicstyle=\ttfamily\footnotesize,
  keywordstyle=\color{StAndrewsBlue}\bfseries,
  stringstyle=\color{StAndrewsMidGreen},
  commentstyle=\color{StAndrewsOrange}\itshape,
  numberstyle=\tiny\color{gray},
  numbers=left,
  numbersep=5pt,
  backgroundcolor=\color{gray!5},
  frame=single,
  rulecolor=\color{StAndrewsBlue!30},
  breaklines=true,
  showstringspaces=false,
  tabsize=4,
  xleftmargin=1.5em,
  framexleftmargin=1em,
  captionpos=b,
  morekeywords={True, False, None, print, range, len, sum, max, min, type,
                exit, int, float, str, bool, import, as, from, for, in,
                while, def, return, class, with, open, f},
  literate={>>>}{{\textcolor{StAndrewsBurgundy}{>{>}>}}}3,
}

% Inline code command
\newcommand{\py}[1]{\lstinline[language=Python]{#1}}

% Shell listing style (for terminal commands)
\lstdefinestyle{shell}{
  language={},
  basicstyle=\ttfamily\footnotesize,
  keywordstyle={},
  stringstyle={},
  commentstyle=\color{StAndrewsOrange}\itshape,
  backgroundcolor=\color{gray!5},
  frame=single,
  rulecolor=\color{StAndrewsBlue!30},
  numbers=none,
  xleftmargin=1.5em,
  framexleftmargin=1em,
  breaklines=true,
  showstringspaces=false,
  literate={~}{{\raise.25ex\hbox{$\sim$}}}1
           {\$}{{\textcolor{StAndrewsMidGreen}{\$}}}1,
}

% Output style (for Python output, light blue border)
\lstdefinestyle{output}{
  language={},
  basicstyle=\ttfamily\scriptsize,
  keywordstyle={},
  stringstyle={},
  commentstyle={},
  backgroundcolor=\color{gray!3},
  frame=single,
  rulecolor=\color{StAndrewsLightBlue!50},
  numbers=none,
  xleftmargin=0.5em,
  framexleftmargin=0.2em,
  breaklines=true,
  showstringspaces=false,
}

% ============================================================================
% ADDITIONAL PACKAGES
% ============================================================================
\usepackage{tcolorbox}

% Custom table column types
\newcolumntype{R}{>{\raggedright\arraybackslash}p{8cm}}

% Highlight box (encouragement / key takeaways)
\newtcolorbox{highlightbox}{
  colback=StAndrewsLightBlue!8,
  colframe=StAndrewsBlue!25,
  boxrule=0.8pt,
  arc=4pt,
  left=6pt, right=6pt,
  top=4pt, bottom=4pt,
}

% Demo box (live demonstration callouts)
\newtcolorbox{demobox}[1][Live Demonstration]{
  colback=StAndrewsBlue!5,
  colframe=StAndrewsBlue!50,
  boxrule=0.8pt,
  arc=4pt,
  left=6pt, right=6pt,
  top=4pt, bottom=4pt,
  fonttitle=\bfseries,
  title=#1,
}


\begin{document}

% ============================================================================
% TITLE PAGE
% ============================================================================
\begin{frame}[plain]
\titlepage
\end{frame}


% ============================================================================
% SECTION 0: OPENING
% ============================================================================
\section{Introduction}
\subsection{Roadmap}

\begin{frame}{What We Will Cover Today}

This is a \alert{hands-on session}. Have your laptop open and follow along.

\medskip

\begin{enumerate}
    \item \textbf{Positron} --- getting a Python session running
    \item \textbf{Console \& Files} --- two ways to write Python
    \item \textbf{\texttt{uv}} --- installing packages into your project
    \item \textbf{Data Structures} --- organising financial data
    \item \textbf{Loops} --- automating repetitive calculations
\end{enumerate}

\medskip

\begin{highlightbox}
\centering
No programming experience required! We will go step by step.
\end{highlightbox}

\end{frame}


\begin{frame}{Table of Contents}
\begin{columns}[t]
  \begin{column}{0.33\textwidth}
    \tableofcontents[sections={1-3}]
  \end{column}
  \begin{column}{0.33\textwidth}
    \tableofcontents[sections={4-6}]
  \end{column}
  \begin{column}{0.33\textwidth}
    \tableofcontents[sections={7-8}]
  \end{column}
\end{columns}
\end{frame}


\subsection{Motivation}
\begin{frame}{Why Python?}

\bi
    \item Used virtually everywhere
    \item Free and open source (unlike Stata, EViews or Bloomberg Terminal)
    \item Huge ecosystem of tools for data, statistics, and visualisation
    \item The most-requested programming skill in finance job postings
    \item You do \textbf{not} need to be a ``programmer'' --- think of it as a very powerful calculator
    \item With AI tools like GitHub Copilot, writing code is easier than ever
\ei

\end{frame}


\subsection{Setup Recap}
\begin{frame}{Recap: What You Already Set Up}

In the \alert{Python Setup Guide} you installed:

\medskip

\begin{tabular}{@{}lp{10cm}@{}}
\toprule
\textbf{Tool} & \textbf{What it does} \\
\midrule
\texttt{uv} & Manages Python and packages \\[0.5em]
Positron & Your code editor (like Word, but for code) \\[0.5em]
GitHub Copilot & AI assistant that helps you write code \\
\bottomrule
\end{tabular}

\bigskip

\begin{highlightbox}
\centering
Today we will learn \textbf{what} these tools actually do and \textbf{how} to use them for finance.
\end{highlightbox}

\end{frame}


% ============================================================================
% GETTING STARTED IN POSITRON
% ============================================================================
\section{Getting Started with Positron}
\subsection{Create a Project Folder}

\begin{frame}{Getting Started in Positron}
\framesubtitle{Step 1: Create an empty folder for your course (e.g.\ FI5699 in Documents)}

\vspace{-0.8em}
\begin{center}
\includegraphics[width=\textwidth,height=\textheight,keepaspectratio]{w2-1-folder.png}
\end{center}

\end{frame}


\begin{frame}{Getting Started in Positron}
\framesubtitle{Step 2: Open Positron and click the Open button}

\vspace{-0.8em}
\begin{center}
\includegraphics[width=\textwidth,height=\textheight,keepaspectratio]{w2-2-positron.png}
\end{center}

\end{frame}


\begin{frame}{Getting Started in Positron}
\framesubtitle{Step 3: Navigate to your folder and click Open}

\vspace{-0.8em}
\begin{center}
\includegraphics[width=\textwidth,height=\textheight,keepaspectratio]{w2-3-positron-open.png}
\end{center}

\end{frame}


\begin{frame}{Getting Started in Positron}
\framesubtitle{Step 4: Your folder is now open --- you'll see the Welcome tab and sidebar}

\vspace{-0.8em}
\begin{center}
\includegraphics[width=\textwidth,height=\textheight,keepaspectratio]{w2-4-positron-folder.png}
\end{center}

\end{frame}


\subsection{Create a Virtual Environment}
\begin{frame}{Getting Started in Positron}
\framesubtitle{Step 5: In the terminal, type \texttt{uv venv --python 3.10} and press Enter}

\vspace{-0.8em}
\begin{center}
\includegraphics[width=\textwidth,height=\textheight,keepaspectratio]{w2-5-positron-uv.png}
\end{center}

\end{frame}


\begin{frame}{Getting Started in Positron}
\framesubtitle{Step 6: Positron detects the new environment --- click Yes}

\vspace{-0.8em}
\begin{center}
\includegraphics[width=\textwidth,height=\textheight,keepaspectratio]{w2-6-positron-uv-yes.png}
\end{center}

\end{frame}


\subsection{Start a Python Session}
\begin{frame}{Getting Started in Positron}
\framesubtitle{Step 7: A Python interpreter appears in the top-right toolbar --- click on it}

\vspace{-0.8em}
\begin{center}
\includegraphics[width=\textwidth,height=\textheight,keepaspectratio]{w2-7-positron-session.png}
\end{center}

\end{frame}


\begin{frame}{Getting Started in Positron}
\framesubtitle{Step 8: Choose New Interpreter Session}

\vspace{-0.8em}
\begin{center}
\includegraphics[width=\textwidth,height=\textheight,keepaspectratio]{w2-8-positron-session.png}
\end{center}

\end{frame}


\begin{frame}{Getting Started in Positron}
\framesubtitle{Step 9: Select the one for your project --- e.g.\ Python 3.10.11 (Uv: FI5699)}

\vspace{-0.8em}
\begin{center}
\includegraphics[width=\textwidth,height=\textheight,keepaspectratio]{w2-9-positron-session.png}
\end{center}

\end{frame}


\begin{frame}{Getting Started in Positron}
\framesubtitle{Step 10: You should see the Python Console with the \texttt{>>>} prompt --- you're ready!}

\vspace{-0.8em}
\begin{center}
\includegraphics[width=\textwidth,height=\textheight,keepaspectratio]{w2-10-positron-session-go.png}
\end{center}

\end{frame}


% ============================================================================
% SECTION 1: THE PYTHON CONSOLE
% ============================================================================
\section{The Python Console}
\subsection{What It Is}

\begin{frame}{The Python Console}

You just opened the Python Console in Positron (Step 10). Let's understand what it is.

\bigskip

\begin{columns}[T]
\begin{column}{0.55\textwidth}
\bi
    \item The \alert{Console} is where you type Python and get instant results
    \item The \texttt{>>>} prompt means Python is waiting for your input
    \item Type something, press Enter, see the answer
    \item Think of it as a very powerful calculator
    \item Positron also has a \alert{Terminal} tab --- that's for system commands (installing packages, etc.)
\ei
\end{column}
\begin{column}{0.42\textwidth}
\centering
\begin{tikzpicture}[node distance=0.8cm, >=stealth, thick,
  box/.style={rectangle, rounded corners=3pt, minimum width=3cm,
              minimum height=0.7cm, align=center, font=\scriptsize}]
  \node[box, draw=StAndrewsBlue!40, fill=StAndrewsBlue!8] (you)
    {You type\\[-1pt]\texttt{2 + 2}};
  \node[box, below=of you, fill=StAndrewsBlue, text=white] (console)
    {Python\\[-1pt]Console};
  \node[box, below=of console, draw=StAndrewsBlue!40, fill=StAndrewsBlue!8] (result)
    {Python replies\\[-1pt]\texttt{4}};
  \draw[->, StAndrewsBlue, line width=0.8pt] (you) -- (console);
  \draw[->, StAndrewsBlue, line width=0.8pt] (console) -- (result);
\end{tikzpicture}
\end{column}
\end{columns}

\end{frame}


\subsection{Console vs Terminal}
\begin{frame}{Console vs Terminal}

Positron has \textbf{two} tabs at the bottom. Make sure you're in the right one:

\bigskip

\begin{tabular}{@{}lll@{}}
\toprule
 & \textbf{Console} & \textbf{Terminal} \\
\midrule
\textbf{Prompt} & \texttt{>>>} & \texttt{\$} (Mac/Linux) or \texttt{PS>} (Windows) \\[0.5em]
\textbf{Language} & Python & Shell (system commands) \\[0.5em]
\textbf{Use for} & Running Python code & Installing packages, file management \\[0.5em]
\textbf{Example} & \texttt{2 + 2} $\to$ \texttt{4} & \texttt{uv add polars} \\
\bottomrule
\end{tabular}

\bigskip

\begin{highlightbox}
\centering
For this session, we will mostly use the \textbf{Console} tab.
\end{highlightbox}

\end{frame}


\subsection{Core Syntax}
\begin{frame}[fragile]{Python as a Calculator}

In the Console, try typing these expressions:

\begin{lstlisting}[style=shell]
>>> 2 + 2
4
>>> 100 * 1.05    # 5% return on 100
105.0
>>> 175000 / 12   # monthly salary
14583.333333333334
\end{lstlisting}

\medskip

\begin{tabular}{@{}lll@{}}
\toprule
\textbf{Operator} & \textbf{Meaning} & \textbf{Example} \\
\midrule
\texttt{+} & Add & \texttt{100 + 50} $\to$ \texttt{150} \\
\texttt{-} & Subtract & \texttt{100 - 30} $\to$ \texttt{70} \\
\texttt{*} & Multiply & \texttt{50 * 3} $\to$ \texttt{150} \\
\texttt{/} & Divide & \texttt{100 / 3} $\to$ \texttt{33.33\ldots} \\
\texttt{**} & Power & \texttt{1.07 ** 10} $\to$ \texttt{1.967\ldots} \\
\bottomrule
\end{tabular}

\end{frame}


\subsection{Variables and Output}
\begin{frame}[fragile]{Storing Values with Variables}

Instead of retyping numbers, give them names:

\begin{lstlisting}
>>> price = 142.50
>>> shares = 100
>>> position = price * shares
>>> position
14250.0
\end{lstlisting}

\medskip

\bi
    \item A \alert{variable} is a name that stores a value
    \item Use \texttt{=} to assign: the name goes on the left, the value on the right
    \item Python remembers it for the rest of your session
    \item You can use variables in calculations, just like in a spreadsheet formula
\ei

\end{frame}


\begin{frame}[fragile]{Displaying Results with \texttt{print()}}

The console shows the result of the last expression automatically. But to display something explicitly, use \py{print()}:

\begin{lstlisting}
>>> ticker = "AAPL"
>>> price = 142.50
>>> shares = 100
>>> print(ticker)
AAPL
>>> print(f"Position value: ${price * shares:,.2f}")
Position value: $14,250.00
\end{lstlisting}

\medskip

\bi
    \item \py{print()} displays whatever you put inside the brackets
    \item An \alert{f-string} (note the \texttt{f} before the quotes) lets you embed variables inside text
    \item \texttt{:,.2f} is optional formatting: commas and 2 decimal places
\ei

\end{frame}


\begin{frame}{Helpful Console Tips}

\begin{tabular}{@{}lp{9cm}@{}}
\toprule
\textbf{Shortcut} & \textbf{What it does} \\
\midrule
\textbf{Up arrow} (\textuparrow) & Recall previous commands --- saves retyping \\[0.5em]
\textbf{Tab} & Autocomplete a variable or function name \\[0.5em]
\textbf{Ctrl + C} & Cancel a running command if something goes wrong \\[0.5em]
\py{type(x)} & Find out what kind of value \texttt{x} is (number, text, etc.) \\[0.5em]
\py{help(x)} & Get documentation for any function \\
\bottomrule
\end{tabular}

\bigskip

\begin{highlightbox}
\centering
You don't need to memorise these. Just remember: \textbf{up arrow} to get back what you just typed.
\end{highlightbox}

\end{frame}


\subsection{Exercises}
\begin{frame}[fragile]{Exercise 1: Explore the Console}

\begin{alertblock}{Try It Yourself (3 minutes)}
\begin{enumerate}
    \item In the Positron \textbf{Console}, compute \texttt{1.07 ** 10} --- what is 7\% compounded over 10 years?
    \item Create a variable \texttt{salary = 175000} and compute the monthly amount (\texttt{salary / 12})
    \item Try \texttt{print(f"Monthly: \${salary / 12:,.2f}")}
    \item Use the \textbf{up arrow} to recall and re-run a previous command
\end{enumerate}
\end{alertblock}

\medskip

\begin{highlightbox}
\centering
Stuck? Raise your hand. There are no silly questions here.
\end{highlightbox}

\end{frame}


% ============================================================================
% SECTION 1b: PYTHON FILES (.py)
% ============================================================================
\section{Python Files (\texttt{.py})}
\subsection{Create a File}

\begin{frame}{Python Files (\texttt{.py})}
\framesubtitle{So far we've typed in the Console. Let's save code in a file instead.}

\vspace{-0.8em}
\begin{center}
\includegraphics[width=\textwidth,height=\textheight,keepaspectratio]{w2-11-python-file.png}
\end{center}

\end{frame}


\begin{frame}{Python Files (\texttt{.py})}
\framesubtitle{Step 1: Click the New button (top-left) and choose New File\ldots}

\vspace{-0.8em}
\begin{center}
\includegraphics[width=\textwidth,height=\textheight,keepaspectratio]{w2-12-python-file.png}
\end{center}

\end{frame}


\begin{frame}{Python Files (\texttt{.py})}
\framesubtitle{Step 2: Select Python File from the list}

\vspace{-0.8em}
\begin{center}
\includegraphics[width=\textwidth,height=\textheight,keepaspectratio]{w2-13-python-file.png}
\end{center}

\end{frame}


\begin{frame}{Python Files (\texttt{.py})}
\framesubtitle{Step 3: A new empty file opens in the editor --- this is where you write your code}

\vspace{-0.8em}
\begin{center}
\includegraphics[width=\textwidth,height=\textheight,keepaspectratio]{w2-14-python-file.png}
\end{center}

\end{frame}


\begin{frame}{Python Files (\texttt{.py})}
\framesubtitle{Step 4: Type your code and save with Cmd+S / Ctrl+S --- give it a \texttt{.py} name}

\vspace{-0.8em}
\begin{center}
\includegraphics[width=\textwidth,height=\textheight,keepaspectratio]{w2-15-python-file.png}
\end{center}

\end{frame}


\begin{frame}{Python Files (\texttt{.py})}
\framesubtitle{Step 5: Click the play button ($\blacktriangleright$) to run --- output appears in the Console below}

\vspace{-0.8em}
\begin{center}
\includegraphics[width=\textwidth,height=\textheight,keepaspectratio]{w2-16-python-file.png}
\end{center}

\end{frame}


\subsection{Run Code}
\begin{frame}{Running Code: Whole File vs Line by Line}

\begin{columns}[T]
\begin{column}{0.48\textwidth}
\textbf{Run the whole file}
\bi
    \item Click the $\blacktriangleright$ button at the top of the editor
    \item Runs every line from top to bottom
    \item Best for finished scripts
\ei
\end{column}
\begin{column}{0.48\textwidth}
\textbf{Run a single line}
\bi
    \item Place your cursor on a line
    \item Press \alert{Cmd+Enter} (Mac) or \alert{Ctrl+Enter} (Windows)
    \item That line runs in the Console
\ei
\end{column}
\end{columns}

\bigskip

\begin{highlightbox}
\centering
\textbf{Important:} When running line by line, Python only knows about lines you've already run. Always run lines \textbf{in order from the top} --- e.g.\ run \texttt{x = 5} before \texttt{print(x)}.
\end{highlightbox}

\end{frame}


% ============================================================================
% SECTION 2: INSTALLING PACKAGES WITH UV
% ============================================================================
\section{Installing Packages with \texttt{uv}}
\subsection{Initialize a Project}

\begin{frame}[fragile]{Setting Up Your Project}

You already have a folder open in Positron with a virtual environment. Now let's turn it into a proper \alert{project} so we can install packages.

\bigskip

In the \textbf{Terminal} tab (not the Console), type:

\begin{lstlisting}[style=shell]
$ uv init
\end{lstlisting}

\medskip

This creates two files in your folder:

\medskip

\begin{tabular}{@{}lp{8cm}@{}}
\toprule
\textbf{File} & \textbf{What it does} \\
\midrule
\texttt{pyproject.toml} & The recipe --- lists which packages your project needs \\[0.3em]
\texttt{uv.lock} & The exact versions --- so your code works the same way every time \\
\bottomrule
\end{tabular}

\bigskip

Think of \texttt{pyproject.toml} like a shopping list. You're about to add items to it.

\end{frame}


\subsection{Add Your Toolkit}
\begin{frame}[fragile]{Your Dissertation Toolkit}

These are the packages we'll use throughout the module:

\medskip

\begin{tabular}{@{}llp{5cm}@{}}
\toprule
\textbf{Package} & \textbf{What it does} & \textbf{Think of it as\ldots} \\
\midrule
\texttt{polars} & Work with tables of data & A supercharged spreadsheet \\[0.3em]
\texttt{plotnine} & Create charts and plots & Excel charts, but better \\[0.3em]
\texttt{pyfixest} & Run regressions & Stata's \texttt{reghdfe}, in Python \\[0.3em]
\texttt{diff-diff} & Difference-in-differences & Causal inference toolkit \\
\bottomrule
\end{tabular}

\bigskip

Install them all in one go (in the \textbf{Terminal} tab):

\begin{lstlisting}[style=shell]
$ uv add polars plotnine pyfixest diff-diff
\end{lstlisting}

\end{frame}


\subsection{Virtual Environments}
\begin{frame}[fragile]{What Is a Virtual Environment?}

\begin{columns}[T]
\begin{column}{0.55\textwidth}
When you ran \texttt{uv venv} earlier, it created a \texttt{.venv} folder inside your project.

\medskip

\bi
    \item This is your project's \alert{private toolbox}
    \item Packages you install go here, not on your whole computer
    \item Different projects can use different packages (or different versions) without clashing
    \item Positron already knows about it (you selected it in step 9)
\ei
\end{column}
\begin{column}{0.42\textwidth}
\centering
\begin{tikzpicture}[node distance=0.6cm, >=stealth, thick,
  box/.style={rectangle, rounded corners=3pt, minimum width=3.2cm,
              minimum height=0.6cm, align=center, font=\scriptsize}]
  \node[box, fill=StAndrewsBlue, text=white] (proj)
    {FI5699/ (your folder)};
  \node[box, below=of proj, draw=StAndrewsBlue!40, fill=StAndrewsBlue!8] (venv)
    {\texttt{.venv/} (private toolbox)};
  \node[box, below=0.4cm of venv, draw=StAndrewsMidBlue!40, fill=StAndrewsMidBlue!8, font=\tiny] (p1)
    {polars};
  \node[box, below=0.2cm of p1, draw=StAndrewsMidBlue!40, fill=StAndrewsMidBlue!8, font=\tiny] (p2)
    {plotnine};
  \node[box, below=0.2cm of p2, draw=StAndrewsMidBlue!40, fill=StAndrewsMidBlue!8, font=\tiny] (p3)
    {pyfixest, diff-diff, \ldots};
  \draw[->, StAndrewsBlue, line width=0.8pt] (proj) -- (venv);
  \draw[StAndrewsMidBlue!50, line width=0.5pt] (venv.south) -- (p1.north);
\end{tikzpicture}
\end{column}
\end{columns}

\end{frame}


\subsection{Common Errors}
\begin{frame}[fragile]{What Happens If You Forget \texttt{uv add}?}

If you try to use a package you haven't installed:

\begin{lstlisting}
>>> import polars as pl
\end{lstlisting}

\begin{lstlisting}[style=output]
ModuleNotFoundError: No module named 'polars'
\end{lstlisting}

\medskip

This is Python telling you: \textit{``I don't have that tool in my toolbox.''}

\bigskip

\textbf{The fix is always the same} --- go to the \textbf{Terminal} tab and run:

\begin{lstlisting}[style=shell]
$ uv add polars
\end{lstlisting}

Then try the import again in the Console. It will work.

\bigskip

\begin{highlightbox}
\centering
\texttt{ModuleNotFoundError} = you need to \texttt{uv add} the package first.
\end{highlightbox}

\end{frame}


\begin{frame}[fragile]{Try It: Import and Use \texttt{polars}}

After running \texttt{uv add polars}, switch to the \textbf{Console} and type:

\begin{lstlisting}
>>> import polars as pl
>>> df = pl.DataFrame({
...     "ticker": ["AAPL", "MSFT", "TSLA"],
...     "price":  [142.50, 378.90, 248.50],
... })
>>> print(df)
\end{lstlisting}

\begin{lstlisting}[style=output]
shape: (3, 2)
+--------+-------+
| ticker | price |
| ---    | ---   |
| str    | f64   |
+========+=======+
| AAPL   | 142.5 |
| MSFT   | 378.9 |
| TSLA   | 248.5 |
+--------+-------+
\end{lstlisting}

You just created a table of data in Python. Don't worry about the \texttt{\{~\}} and \texttt{[~]} syntax --- we'll cover those after the break.

\end{frame}


\begin{frame}[fragile]{The Three \texttt{uv} Commands You Need}

\begin{tabular}{@{}lp{7.5cm}@{}}
\toprule
\textbf{What you want to do} & \textbf{Command (in the Terminal tab)} \\
\midrule
Set up a new project & \texttt{uv init} \\[0.5em]
Install a package & \texttt{uv add polars} \\[0.5em]
See what's installed & \texttt{uv pip list} \\
\bottomrule
\end{tabular}

\bigskip

That's it. Three commands. Everything else is just Python.

\bigskip

\begin{highlightbox}
\centering
\textbf{The pattern:} \texttt{uv add \ldots} in the Terminal, then \texttt{import \ldots} in the Console (or your \texttt{.py} file).
\end{highlightbox}

\end{frame}


\subsection{Exercises}
\begin{frame}[fragile]{Exercise 2: Install and Import}

\begin{alertblock}{Try It Yourself (3 minutes)}
\begin{enumerate}
    \item In the \textbf{Terminal} tab, run: \texttt{uv add polars plotnine pyfixest diff-diff}
    \item Watch it install --- it should only take a few seconds
    \item In the \textbf{Console}, type: \texttt{import polars as pl}
    \item Then type: \texttt{print(pl.\_\_version\_\_)} --- you should see a version number
    \item Open \texttt{pyproject.toml} in the editor --- find where \texttt{polars} appears
\end{enumerate}
\end{alertblock}

\medskip

\begin{highlightbox}
\centering
If you see \texttt{ModuleNotFoundError}, check: are you in the \textbf{Console} (not Terminal)? Did \texttt{uv add} finish?
\end{highlightbox}

\end{frame}


% ============================================================================
% SECTION 3: DATA STRUCTURES
% ============================================================================
\section{Data Structures}
\subsection{Why It Matters}

\begin{frame}{Why Data Structures Matter}

Later in the module, \texttt{polars} will handle your actual datasets (tables of stock prices, returns, etc.). So why learn lists and dictionaries?

\bigskip

Because they are the \alert{building blocks} that show up \textbf{everywhere} in Python:

\medskip

\begin{tabular}{@{}lp{7.5cm}@{}}
\toprule
\textbf{Situation} & \textbf{Example} \\
\midrule
Passing options to a function & \texttt{yf.download(["AAPL", "MSFT"], \ldots)} \\[0.3em]
Storing settings or parameters & \texttt{\{"start": "2024-01-01", "end": "2024-12-31"\}} \\[0.3em]
Collecting results in a loop & \texttt{rows.append(\{"ticker": t, "return": r\})} \\[0.3em]
Organising any small group of values & \texttt{rates = [0.05, 0.07, 0.10]} \\
\bottomrule
\end{tabular}

\bigskip

\begin{highlightbox}
\centering
Think of lists and dictionaries as Python's \textbf{plumbing} --- you won't store large datasets in them, but you can't write a useful program without them.
\end{highlightbox}

\end{frame}


\subsection{Variables and Types}
\begin{frame}[fragile]{Variables: Giving Names to Values}

\begin{lstlisting}
price = 142.50          # a decimal number
shares = 100            # a whole number
ticker = "AAPL"         # text (always in quotes)
is_profitable = True    # yes or no (True / False)

print(f"Ticker: {ticker}")
print(f"Price:  {price}")
print(f"Value:  {price * shares}")
\end{lstlisting}

\begin{lstlisting}[style=output]
Ticker: AAPL
Price:  142.5
Value:  14250.0
\end{lstlisting}

\begin{highlightbox}
\centering
A variable is just a \textbf{name tag} you stick on a value. Python figures out the type automatically.
\end{highlightbox}

\end{frame}


\subsection{Lists}
\begin{frame}[fragile]{Lists: Ordered Collections}

\begin{lstlisting}
# A portfolio of stock tickers
portfolio = ["AAPL", "MSFT", "TSLA", "AMZN"]

print(portfolio[0])     # first item  (counting starts at 0!)
print(portfolio[-1])    # last item
print(len(portfolio))   # how many items
\end{lstlisting}

\begin{lstlisting}[style=output]
'AAPL'
'AMZN'
4
\end{lstlisting}

\bi
    \item Created with square brackets \texttt{[~]}
    \item Items have a fixed \alert{order} (first, second, third\ldots)
    \item Counting starts at \textbf{0}, not 1 (this trips everyone up at first!)
\ei

\end{frame}


\begin{frame}[fragile]{Useful Things You Can Do with Lists}

\begin{lstlisting}
portfolio = ["AAPL", "MSFT"]
portfolio.append("GOOGL")       # add to the end
portfolio.remove("MSFT")        # remove by name
print(portfolio)                # ['AAPL', 'GOOGL']

returns = [0.5, -1.2, 0.8, 2.1, -0.3]
print(sum(returns))      # total:  1.9
print(max(returns))      # best:   2.1
print(min(returns))      # worst: -1.2
print(len(returns))      # count:  5
\end{lstlisting}

\medskip

Built-in functions like \py{sum()}, \py{max()}, \py{min()}, and \py{len()} work directly on lists.

\end{frame}


\begin{frame}{Visualising Those Returns}

\begin{center}
\includegraphics[width=0.8\textwidth]{returns_bar.pdf}
\end{center}

\centering
Blue bars = positive days. Red bars = negative days. (We'll learn to make charts like this next week with \texttt{plotnine}.)

\end{frame}


\subsection{Dictionaries}
\begin{frame}[fragile]{Dictionaries: Labelled Data}

\begin{lstlisting}
# Stock info as a dictionary
stock = {
    "ticker": "AAPL",
    "price": 142.50,
    "shares": 100,
    "sector": "Technology",
}

print(stock["price"])          # look up by name
stock["price"] = 145.00        # update a value
print(stock["price"])
\end{lstlisting}

\bi
    \item Created with curly braces \texttt{\{~\}}
    \item Each entry has a \alert{name} (key) and a \alert{value}
    \item Like a row in a spreadsheet where each column has a header
\ei

\end{frame}


\begin{frame}{Lists vs Dictionaries}

\begin{tabular}{@{}lp{4.5cm}p{5.5cm}@{}}
\toprule
 & \textbf{List} & \textbf{Dictionary} \\
\midrule
\textbf{Looks like} & \texttt{["AAPL", "MSFT"]} & \texttt{\{"ticker": "AAPL", "price": 142\}} \\[0.5em]
\textbf{Access by} & Position (number) & Name (key) \\[0.5em]
\textbf{Best for} & Sequences of similar items & A record with named fields \\[0.5em]
\textbf{Analogy} & A column in Excel & A row in Excel \\
\bottomrule
\end{tabular}

\bigskip

For this course, these two are the ones you'll use most. We'll mention \textbf{tuples} \texttt{(1, 2, 3)} and \textbf{sets} \texttt{\{1, 2, 3\}} as we go, but don't worry about them now.

\end{frame}


\subsection{Putting It Together}
\begin{frame}[fragile]{Live: Building a Portfolio}

\begin{demobox}
Create a dictionary of stock positions and compute values.
\end{demobox}

\begin{lstlisting}
portfolio = {
    "AAPL": {"shares": 50, "price": 142.50},
    "MSFT": {"shares": 30, "price": 378.90},
}

# Total value of the Apple position
aapl_value = (portfolio["AAPL"]["shares"]
              * portfolio["AAPL"]["price"])
print(f"AAPL position value: ${aapl_value:,.2f}")
\end{lstlisting}

\begin{lstlisting}[style=output]
AAPL position value: $7,125.00
\end{lstlisting}

\end{frame}


\subsection{Exercises}
\begin{frame}[fragile]{Exercise 3: Data Structures}

\begin{alertblock}{Try It Yourself (5 minutes)}
\begin{enumerate}
    \item Create a \textbf{list} of 5 stock tickers you know (e.g.\ \texttt{["AAPL", "TSLA", \ldots]})
    \item Print the first and last ticker
    \item Create a \textbf{dictionary} for one stock with keys \texttt{"ticker"}, \texttt{"price"}, \texttt{"shares"}
    \item Print the total value (price $\times$ shares)
\end{enumerate}
\end{alertblock}

\medskip

\textbf{Hint:} Create a new Python file \texttt{exercise3.py} (New $\to$ New File $\to$ Python File), write your code, and click $\blacktriangleright$ to run it.

\end{frame}


% ============================================================================
% SECTION 4: LOOPS
% ============================================================================
\section{Loops}
\subsection{Motivation}

\begin{frame}{Why Do We Need Loops?}

\begin{columns}[T]
\begin{column}{0.52\textwidth}
Computing the return for \textbf{1 stock} is easy:

\medskip

\texttt{value = 100 * 142.50}

\medskip

But what about \textbf{500 stocks}?

\medskip

A \alert{loop} says: \textit{``do this for \textbf{each} item in a list''}

\medskip

It's like dragging a formula down a column in Excel --- but more powerful.
\end{column}
\begin{column}{0.45\textwidth}
\centering
\begin{tikzpicture}[node distance=1cm, >=stealth, thick,
  box/.style={rectangle, rounded corners=3pt, minimum width=2.8cm,
              minimum height=0.7cm, align=center, font=\scriptsize},
  decision/.style={diamond, aspect=2, minimum width=2cm,
                   align=center, font=\scriptsize}]
  \node[box, fill=StAndrewsBlue, text=white] (start)
    {Start: list of\\[-1pt]tickers};
  \node[decision, below=of start, draw=StAndrewsMidBlue,
        fill=StAndrewsMidBlue!15] (check)
    {More\\tickers?};
  \node[box, right=1.8cm of check, fill=StAndrewsMidBlue, text=white] (process)
    {Process\\[-1pt]next ticker};
  \node[box, below=1.2cm of check, fill=StAndrewsGreen, text=white] (done)
    {Done!};
  \draw[->, line width=0.8pt] (start) -- (check);
  \draw[->, line width=0.8pt] (check) -- node[above, font=\tiny] {Yes} (process);
  \draw[->, line width=0.8pt] (check) -- node[right, font=\tiny] {No} (done);
  \draw[->, line width=0.8pt, rounded corners=3pt]
    (process.north) -- ++(0,0.5) -| (check);
\end{tikzpicture}
\end{column}
\end{columns}

\end{frame}


\subsection{\texttt{for} Loops}
\begin{frame}[fragile]{Your First \texttt{for} Loop}

\begin{lstlisting}
tickers = ["AAPL", "MSFT", "TSLA", "AMZN"]

for ticker in tickers:
    print(f"Analysing {ticker}...")
\end{lstlisting}

\medskip

\lstinputlisting[style=output]{outputs/for_loop.txt}

\medskip

\begin{highlightbox}
\centering
\textbf{How to read it:} ``For each ticker in the list, print a message.''\\
The indented line (4 spaces) is what repeats. Everything else runs once.
\end{highlightbox}

\end{frame}


\subsection{Counting with \texttt{range()}}
\begin{frame}[fragile]{\texttt{range()}: Counting by Numbers}

\begin{lstlisting}
# Compound interest: 10 years at 5%
value = 1000
for year in range(1, 11):
    value = value * 1.05
    print(f"Year {year:2d}: ${value:>10,.2f}")
\end{lstlisting}

\medskip

\lstinputlisting[style=output]{outputs/compound.txt}

\medskip

\bi
    \item \py{range(5)} gives you 0, 1, 2, 3, 4
    \item \py{range(1, 11)} gives you 1, 2, 3, \ldots, 10
    \item Useful for simulating time periods
\ei

\end{frame}


\subsection{Looping with Indexes}
\begin{frame}[fragile]{Loops with Calculations}

\begin{lstlisting}
prices = [142.50, 378.90, 248.50, 186.75]
shares = [50, 30, 20, 40]

total = 0
for i in range(len(prices)):
    position = prices[i] * shares[i]
    total += position              # shorthand for total = total + position
    print(f"Position {i+1}: ${position:,.2f}")

print(f"\nPortfolio value: ${total:,.2f}")
\end{lstlisting}

\medskip

\lstinputlisting[style=output]{outputs/portfolio_loop.txt}

\end{frame}


\begin{frame}{Compound Interest: Visualised}

\begin{center}
\includegraphics[width=\textwidth]{compound_growth.pdf}
\end{center}

\smallskip

Same starting amount, different rates. Small differences compound into huge gaps over time.

\end{frame}


\subsection{Nested Loops}
\begin{frame}[fragile]{The Code Behind That Chart}

\begin{lstlisting}
principal = 10000
rates = {"5%": 0.05, "7%": 0.07, "10%": 0.10}

rows = []
for label, rate in rates.items():      # .items() gives each key-value pair
    value = principal
    for year in range(0, 21):           #   for each year...
        rows.append({"rate": label,
                     "year": year,
                     "value": value})
        value = value * (1 + rate)      #   grow by the rate

# This is a nested loop: a loop inside a loop!
\end{lstlisting}

\medskip

Two loops working together: the outer loop picks a rate, the inner loop simulates 20 years.

\end{frame}


\subsection{\texttt{while} Loops}
\begin{frame}[fragile]{\texttt{while} Loops: Repeat Until a Condition}

\begin{lstlisting}
# How many years to double your money at 7%?
value = 1000
target = 2000
years = 0

while value < target:
    value = value * 1.07
    years += 1                # add 1 to years

print(f"It takes {years} years to double at 7%")
\end{lstlisting}

\begin{lstlisting}[style=output]
It takes 11 years to double at 7%
\end{lstlisting}

\begin{highlightbox}
\centering
Use \texttt{for} when you know how many times to repeat. Use \texttt{while} when you want to keep going \textbf{until} something happens.
\end{highlightbox}

\end{frame}


\subsection{Combining with Data Structures}
\begin{frame}[fragile]{Combining Loops and Data Structures}

\begin{columns}[T]
\begin{column}{0.52\textwidth}
\begin{lstlisting}[basicstyle=\ttfamily\scriptsize]
portfolio = {
    "AAPL": {"shares": 50,
             "price": 142.50},
    "MSFT": {"shares": 30,
             "price": 378.90},
    "TSLA": {"shares": 20,
             "price": 248.50},
}

total = 0
for ticker, info in portfolio.items():  # each key + value
    value = info["shares"] * info["price"]
    total += value
    print(f"{ticker}: ${value:,.2f}")

print(f"\nTotal: ${total:,.2f}")
\end{lstlisting}
\end{column}
\begin{column}{0.45\textwidth}
\centering
\includegraphics[width=\textwidth]{portfolio_bar.pdf}

\bigskip

{\small Loop over structured data to compute everything at once.}
\end{column}
\end{columns}

\end{frame}


\subsection{Exercises}
\begin{frame}[fragile]{Exercise 4: Loops}

\begin{alertblock}{Try It Yourself (5 minutes)}
\begin{enumerate}
    \item Write a \py{for} loop that prints each ticker in your list from Exercise 3
    \item Compound returns exercise:
    \begin{itemize}
        \item Start with \texttt{value = 1.0}
        \item For each return in \texttt{[0.05, -0.02, 0.03, 0.01, -0.01]}, do: \texttt{value = value * (1 + r)}
        \item Print the final value
    \end{itemize}
    \item \textbf{Bonus:} Use a \py{while} loop to find how many years \$5,000 takes to reach \$10,000 at 6\%
\end{enumerate}
\end{alertblock}

\end{frame}


% ============================================================================
% SECTION 5: SUMMARY
% ============================================================================
\section{Summary}
\subsection{Key Takeaways}

\begin{frame}{What You Learned Today}

\begin{columns}[T]
\begin{column}{0.48\textwidth}
\textbf{Tools \& setup}
\begin{itemize}
    \item Positron: Console vs Terminal vs Editor
    \item Python files (\texttt{.py}) and how to run them
    \item \texttt{uv init} / \texttt{uv add} for packages
    \item Virtual environments (your project's private toolbox)
\end{itemize}
\end{column}
\begin{column}{0.48\textwidth}
\textbf{Python fundamentals}
\begin{itemize}
    \item Variables and types (\py{int}, \py{float}, \py{str}, \py{bool})
    \item Lists \texttt{[~]} and dictionaries \texttt{\{~\}}
    \item \py{for} loops and \py{while} loops
    \item \py{range()} for counting
    \item Combining data structures with loops
\end{itemize}
\end{column}
\end{columns}

\end{frame}


\subsection{Self-Check}
\begin{frame}[fragile]{Quick Self-Check}

\textbf{Can you answer these in your own words?}

\be
    \item What is the difference between the Console and the Terminal in Positron?
    \item What does \texttt{uv add polars} do, and where do you run it?
    \item What is the difference between a list and a dictionary?
    \item What does \py{for ticker in portfolio:} do?
    \item When would you use \py{while} instead of \py{for}?
\ee

\end{frame}


\subsection{Next Steps}
\begin{frame}{What's Next}

\bi
    \item Week 4: \alert{real financial data} with \texttt{polars} and \texttt{plotnine}
    \begin{itemize}
        \item Downloading stock prices, filtering tables, computing returns, making charts
    \end{itemize}
    \item \textbf{Homework:} complete today's exercises and experiment in the Console
    \item Resources:
    \begin{itemize}
        \item Python tutorial: \url{https://docs.python.org/3/tutorial/}
        \item \texttt{uv} docs: \url{https://docs.astral.sh/uv/}
        \item polars guide: \url{https://docs.pola.rs/}
        \item plotnine docs: \url{https://plotnine.org/}
    \end{itemize}
\ei

\medskip

\begin{highlightbox}
\centering
If something breaks: read the error message, restart your terminal, try again.\\
Still stuck? Email me at \href{mailto:ce50@st-andrews.ac.uk}{ce50@st-andrews.ac.uk}.
\end{highlightbox}

\end{frame}


\end{document}
